\item[\textbf{Problem 5}]
Suppose the sub-key generation function is to reverse all the bits of the key
$K$ (e.g., $0111 \rightarrow 1110$), and the scrambling function is $f = M
{\XOR} K'$, where $M$ is the second half of input bits and $K'$ is the
sub-key (i.e., the reverse of the original key $K$). Now given the original
$K=0011$, and input bits $1111\ 0000$, compute the output of the Feistel Cipher.

\Solution

Recall, the Feistel cipher looks as follows:

\begin{minipage}{\linewidth}
    \centering
    \includegraphics[width=0.8\linewidth]{feistel-cipher.pdf}
    \captionof{figure}{Feistel Cipher Diagram}
\end{minipage}

We have the following inputs:
\begin{alignat*}{2}
    \text{Input} &= 1111\ 0000                 \\
    K            &= 0011                       \\
    f            &= M \XOR K'                  \\
    M            &= \text{Right}(\text{Input}) \\
\end{alignat*}

Based on the problem description, we have:
\begin{alignat*}{2}
    K'    &= 1100 \\
    M     &= 0000 \\
    R_{0} &= M    \\
    L_{0} &= 1111 \\
\end{alignat*}

Therefore, we can draw the following diagram:

\begin{minipage}{\linewidth}
    \centering
    \includegraphics[width=0.8\linewidth]{problem-5-step-1.pdf}
    \captionof{figure}{Substituting into $F_{k_1}$}
\end{minipage}

And here is a step by step solution:

\begin{minipage}{\linewidth}
    \centering
    \includegraphics[width=0.8\linewidth]{problem-5-step-2.pdf}
    \captionof{figure}{Propagating values for $L_1$ and $R_1$ by following the diagram}
\end{minipage}

\begin{minipage}{\linewidth}
    \centering
    \includegraphics[width=0.8\linewidth]{problem-5-step-3.pdf}
    \captionof{figure}{Plugging in the values for $L_0$, $R_0$, and $F_{k_1}$}
\end{minipage}

\begin{minipage}{\linewidth}
    \centering
    \includegraphics[width=0.8\linewidth]{problem-5-step-4.pdf}
    \captionof{figure}{SimpSimplifying $F_{k_1}$, propagating $L_1$, and simplifying the calculation for $R_1$}
\end{minipage}

\begin{minipage}{\linewidth}
    \centering
    \includegraphics[width=0.8\linewidth]{problem-5-step-5.pdf}
    \captionof{figure}{Simplifying the result for $R_1$}
\end{minipage}

Here is the final answer:
\[
    \text{Output} = 0000\ 0011
\]
